\section*{第 \ref{ch:b3:spectral} 章 --- 紧算子
与谱定理}

\begin{solution}{exo:b3:spectral:1}
(a) $T(B)$ 是有限维空间
$\operatorname{im}T$ 的有界子集，由海涅–博雷尔定理知其相对紧
（\cref{cor:b3:topology:heineborel}，经
与 $\R^n$ 的线性同胚传递）。
(b) $I$ 紧，即闭单位球紧，由
里斯定理（本科二年级）可知，此恰在有限
维时成立。若紧算子 $T$ 有有界逆 $T^{-1}$，则
$I = T^{-1}T$ 将是紧的
（\cref{prop:b3:spectral:ideal}）：在无限
维中不可能。
\end{solution}

\begin{solution}{exo:b3:spectral:2}
(a) $\norm{Tx}^2 = \sum\abs{\lambda_n}^2\abs{x_n}^2 \leq
\sup\abs{\lambda_n}^2\norm x^2$，且近似取等号的情形出现在
达到上确界的 $e_n$ 上。
(b) ($\Leftarrow$) 截断 $T_N$（保留 $n \leq N$，
其余置零）秩有限，且 $\vertiii{T - T_N} =
\sup_{n>N}\abs{\lambda_n} \to 0$，由
\cref{prop:b3:spectral:ideal}知其紧。（$\Rightarrow$）若
$\abs{\lambda_{n_k}} \geq \delta > 0$ 沿某子列成立：
$\norm{Te_{n_k} - Te_{n_l}}^2 = \abs{\lambda_{n_k}}^2 +
\abs{\lambda_{n_l}}^2 \geq 2\delta^2$，因此
$(Te_{n_k})$ 无收敛子列。
(c) $T^* = $ 是对角元为 $(\bar\lambda_n)$ 的对角算子，自伴当且仅当
所有 $\lambda_n \in \R$。此时标准基 $(e_n)$ 为
由特征向量构成的标准正交基，特征值 $\lambda_n \to
0$，即谱定理原样适用。
\end{solution}

\begin{solution}{exo:b3:spectral:3}
设 $(x_n)$ 有界，则 $(Bx_n)$ 有界
（$\vertiii B < \infty$）；由 $S$ 的紧性可取出
子列使 $SBx_{n_k} \to y$；由 $A$ 的连续性得 $ASBx_{n_k} \to
Ay$：$ASB$ 紧。若 $ST = TS = I$，且 $T$ 紧，
$\dim H = \infty$ 时，$I = ST$ 将为紧，与
\cref{exo:b3:spectral:1}(b) 相矛盾 ---
此即从另一侧看到的同一陈述。
\end{solution}

\begin{solution}{exo:b3:spectral:4}
(a) 由 $y$ 上的柯西–施瓦茨不等式，$\abs{T_kf(x)}^2 \leq
\bigl(\int\abs{k(x,y)}^2\dd y\bigr)\norm f_2^2$；对
$x$ 积分（托内利定理）：$\norm{T_kf}_2 \leq \norm k_{L^2(\square)}
\norm f_2$。

(b) 函数族 $e_{mn}(x,y) = e_m(x)\overline{e_n(y)}$ 为
$L^2(\intcc01^2)$ 中标准正交（托内利定理可分离二重
积分）。完备性：若 $h \perp$ 所有 $e_{mn}$，则对每个
$m$，函数 $y \mapsto \int h(x,y)\overline{e_m(x)}\dd
x$（由柯西–施瓦茨不等式与托内利定理知其属于 $L^2$）与
每个 $\overline{e_n}$ 正交 --- 而共轭族
$(\overline{e_n})$ 在 $(e_n)$ 为希尔伯特基时亦构成希尔伯特基
（共轭是 $L^2$ 上保持正交性与完备性的
等距双射）--- 故其几乎处处为 $0$；于是对几乎处处的 $y$，$h(\cdot, y) \perp$ 每个
$e_m$：$h(\cdot, y) = 0$ 几乎处处，由托内利定理 $h = 0$。于是
$(e_{mn})$ 为希尔伯特基；将 $k = \sum c_{mn}e_{mn}$ 展开。
截断 $k_N$（指标 $\leq N$）给出的 $T_{k_N}$ 的
有限秩（值域包含于 $\operatorname{Vect}(e_1, \dots,
e_N)$），且由 (a)，
\[
\vertiii{T_k - T_{k_N}} \leq \norm{k - k_N}_{L^2} \to 0 :
\]
$T_k$ 是有限秩算子的范数极限，故紧
（由\cref{prop:b3:spectral:ideal}）。
\end{solution}

\begin{solution}{exo:b3:spectral:5}
(a) $\lambda\norm x^2 = \langle x, Tx\rangle \geq 0$ 对一个
特征向量成立。该公式见
\cref{prop:b3:spectral:sanorm}，所有取值 $\langle x,
Tx\rangle \geq 0$，因此绝对值符号是多余的。
(b) $(x, y) \mapsto \langle x, Ty\rangle$ 为埃尔米特
半正定（允许退化）的半双线性形式；通常的
柯西–施瓦茨证明（展开 $\langle x + t\eu^{\iu\theta}y,
T(x + t\eu^{\iu\theta}y)\rangle \geq 0$，并取
判别式）从未用到正定性。
\end{solution}

\begin{solution}{exo:b3:spectral:6}
(a) $\norm{Mf}_2 \leq \norm f_2$，且在 $f_n = \sqrt n\,
\mathbf 1_{\intcc{1 - 1/n}1}$（均为单位向量），$\norm{Mf_n}_2
\geq 1 - \frac1n$，故 $\vertiii M = 1$；且自伴，因
乘子为实。特征值：若 $xf(x) = \lambda f(x)$ 几乎处处
成立，则 $f = 0$ 在零测集 $\{x = \lambda\}$ 之外几乎处处，即 $f =
0$ 在 $L^2$ 中。
(b) 取同样的 $f_n$：$\norm{Mf_n - f_n}_2 \leq \frac1n \to
0$，而 $f_n \rightharpoonup 0$（对固定的 $g \in L^2$，
$\abs{\langle g, f_n\rangle} \leq \norm{g\,\mathbf
1_{\intcc{1-1/n}1}}_2 \to 0$（由控制收敛定理）。若 $Mf_{n_k} \to h$ 依
范数收敛，则 $f_{n_k} \to h$，迫使 $h = 0$（弱极限），
但 $\norm h = 1$：无收敛子列。
(c) 在\cref{lem:b3:spectral:existence}中，恰好
取出“$Tx_{n_k} \to y$”这一步用到紧性；对 $M$ 而言，
极大化序列集中在 $x = 1$ 附近，其
像弱收敛于 $0$，从不依范数收敛，
故数值域顶端的特征向量根本不存在。
\end{solution}

\begin{solution}{exo:b3:spectral:7}
(a) 取 $V = T_k$，其中 $k(x, y) = \mathbf 1_{y < x} \in
L^2(\intcc01^2)$，紧（见\cref{exo:b3:spectral:4}）。其
伴随为核 $\overline{k(y,
x)} = \mathbf 1_{y > x}$ 的核算子：$V^*f(x) = \int_x^1f$；$V \neq V^*$
（取 $f = \mathbf 1$ 检验）。
(b) 若 $Vf = \lambda f$，$\lambda \ne 0$，则 $Vf$ 连续
于 $\intcc01$ 上（对 $\int_0^x f$ 用控制收敛定理），故 $f
= \frac1\lambda Vf$ 有连续代表元；于是
$Vf$ 为 $\mathcal C^1$（由微积分基本定理，对
连续被积函数），故 $f$ 为 $\mathcal C^1$，且
$\lambda f' = f$，且 $f(0) = \frac1\lambda Vf(0) = 0$：$f =
C\eu^{x/\lambda}$，其中 $C = f(0) = 0$。
(c) $V$ 为紧算子，它除可能的
$0$ 外无特征值（$Vf = 0$ 经对积分求导，推出 $f = 0$ 几乎处处
成立，故连 $0$ 也不是特征值）：谱
理论真正需要自伴性，而不仅是紧性。
\end{solution}

\begin{solution}{exo:b3:spectral:8}
设 $x = \sum_ic_ie_i + z$，$z \in \ker T$，则 $\langle x,
Tx\rangle = \sum_i\mu_i\abs{c_i}^2$。\emph{下界}：
在 $V_k = \operatorname{Vect}(e_1, \dots,
e_k)$ 的单位球面上，$\langle x, Tx\rangle = \sum_{i\leq
k}\mu_i\abs{c_i}^2 \geq \mu_k$，故对 $V$ 先取最小值、再取最大值，
所得值 $\geq \mu_k$。\emph{上界}：设 $\dim V = k$，且 $W_k
= \overline{\operatorname{Vect}}(e_k, e_{k+1}, \dots) +
\ker T$，余维数为 $k - 1$（其正交补为
$V_{k-1}$）；$V \cap W_k \neq \{0\}$（线性映射 $V \to
H/W_k \cong V_{k-1}$，秩 $\leq k - 1$，必有非平凡
核），且单位向量 $x \in V\cap W_k$ 满足 $\langle x,
Tx\rangle = \sum_{i \geq k}\mu_i\abs{c_i}^2 \leq \mu_k$，
故 $V$ 上的最小值 $\leq \mu_k$。综上即得第一式；
第二式可类似证明（取 $W = W_k$ 检验；对
任意 $W$，其余维数为 $k-1$，由 $W \cap V_k \neq 0$，
可得一个单位向量，使 $\langle x, Tx\rangle \geq \mu_k$）。
单调性：$\langle x, Tx\rangle \leq \langle x,
Sx\rangle$ 逐点成立，经 $\max\min$ 传递。
\end{solution}

\begin{solution}{exo:b3:spectral:9}
改写为 $f - \lambda Tf = g$。当 $\lambda = 0$ 时：$f = g$，
总有唯一解。当 $\lambda \neq 0$ 时：此为
$(T - \frac1\lambda)f = -\frac g\lambda$，由弗雷德霍姆择一定理
（见\cref{thm:b3:spectral:fredholm}），并以
$\lambda_n = \bigl((n + \frac12)\pi\bigr)^{-2}$ 为
$T$ 的特征值（见\cref{ex:b3:spectral:minxy}）：对所有 $g$ 唯一可解，当且仅当
$\frac1\lambda \neq \lambda_n$ 对每个 $n$ 成立，即
\[
\lambda \neq \Bigl(n + \tfrac12\Bigr)^2\pi^2
\qquad (n = 0, 1, 2, \dots).
\]
在例外值 $\lambda = (n+\frac12)^2\pi^2$ 处，解
存在当且仅当 $g \perp \sin\bigl((n{+}\frac12)\pi x\bigr)$，且
此时解唯一，至多相差该正弦的任意倍数。
\end{solution}

\begin{solution}{exo:b3:spectral:10}
(a) $Sx = \lambda x$：比较分量，得 $0 = \lambda
x_1$，且 $x_n = \lambda x_{n+1}$；若 $\lambda \ne 0$，则
$x_1 = 0$，归纳可得 $x = 0$；若 $\lambda = 0$，则 $Sx = 0$
推出 $x = 0$（$S$ 为等距算子）。故无特征值。$S^*x =
\lambda x$ 即 $x_{n+1} = \lambda x_n$：$x =
x_1(1, \lambda, \lambda^2, \dots)$，它在 $\ell^2$ 中当且仅当
$\abs\lambda < 1$，故特征值充满一个开圆盘。
(b) $\norm{Se_n - Se_m} = \norm{e_{n+1} - e_{m+1}} = \sqrt2$：
有界列 $(e_n)$ 的像不含柯西子列；
同理 $S^*e_{n+1} = e_n$。两者都不是紧算子。
(c) 对紧自伴算子，特征值
已包含全部内容（见\cref{thm:b3:spectral:spectral}）；去掉
其中任一假设，特征值可能完全消失（$S$，
沃尔泰拉）或充满一个圆盘（$S^*$）：稳健的对象是谱
$\{\lambda : T - \lambda I \text{ 不可逆}\}$，
其理论属于后续课程。
\end{solution}

\begin{solution}{exo:b3:spectral:11}
(a) $S$ 是对角算子，其系数为
$\sqrt{\mu_n} \to 0$，故为紧算子
（由\cref{exo:b3:spectral:2}(b)，转移到
$(e_n)$ 与 $\ker T$ 补全的基上，其中 $S = 0$），自伴
（实对角），正算子（$\langle x, Sx\rangle =
\sum\sqrt{\mu_n}\abs{\langle e_n, x\rangle}^2$），且 $S^2 =
T$ 逐项成立。

(b) $R$ 与 $T = R^2$ 可交换。设 $x$ 是
$T$ 的特征向量，特征值为 $\mu$：$T(Rx) = RTx = \mu Rx$，故
有限维特征子空间 $E_\mu = \ker(T - \mu)$ 关于
$R$ 不变。在 $E_\mu$ 上，$R$ 为对称正算子，且 $R^2
= \mu\,\mathrm{id}$：其特征值 $\rho$ 满足 $\rho^2
= \mu$，$\rho \geq 0$：全部等于 $\sqrt\mu$，而
可对角化且仅有一个特征值的算子是数乘算子：
$R = \sqrt\mu\,\mathrm{id}$ 在 $E_\mu$ 上成立。在 $\ker T$ 上，$\norm
{Rx}^2 = \langle x, R^2x\rangle = \langle x, Tx\rangle = 0$。
故 $R$ 与 $S$ 在 $\ker T$ 及每个特征子空间上一致，
它们的闭线性包为 $H$（由谱定理），故 $R = S$。

(c) $\sqrt G$ 以 $\frac1{n\pi}$ 作用于 $e_n =
\sqrt2\sin(n\pi x)$：它是核算子，其核为
\[
k(x, y) = \sum_{n\geq1}\frac{2\sin(n\pi x)\sin(n\pi
y)}{n\pi},
\]
该级数在 $L^2(\intcc01^2)$ 中收敛（系数
$\frac1{n\pi} \in \ell^2$；部分和核给出
有限秩逼近）。无需初等闭形式，
所需的只是：谱的一侧\emph{就是}算子本身。
\end{solution}

\begin{solution}{exo:b3:spectral:12}
(a) $T^*T$ 是紧的（有界算子与紧算子之积，
见\cref{exo:b3:spectral:3}），自伴
（$(T^*T)^* = T^*T$），正算子（$\langle x, T^*Tx\rangle =
\norm{Tx}^2$）。核：$T^*Tx = 0 \Rightarrow \norm{Tx}^2 =
\langle x, T^*Tx\rangle = 0 \Rightarrow Tx = 0$，且
反之亦然：$\ker T^*T = \ker T$。由谱定理
给出标准正交系 $(e_n)$，使 $T^*Te_n = s_n^2e_n$，
$s_n > 0$，且张成 $(\ker T)^\perp$。

(b) $\langle f_m, f_n\rangle = \frac{\langle Te_m,
Te_n\rangle}{s_ms_n} = \frac{\langle e_m,
T^*Te_n\rangle}{s_ms_n} = \frac{s_n^2}{s_ms_n}\delta_{mn} =
\delta_{mn}$。将 $x = x_0 + \sum_n\langle e_n, x\rangle
e_n$，其中 $x_0 \in \ker T$（在闭线性包与核的和上应用帕塞瓦尔恒等式）；
对连续算子 $T$ 作用，得：
\[
Tx = \sum_n\langle e_n, x\rangle\,Te_n =
\sum_ns_n\langle e_n, x\rangle\,f_n,
\]
该级数收敛，因其部分和为柯西列
（$\norm{\sum_{N<n\leq M}s_n\langle e_n, x\rangle f_n}^2 =
\sum s_n^2\abs{\langle e_n, x\rangle}^2$，
被 $\sup_{n>N}s_n^2\cdot\norm x^2$ 控制，且 $s_n \to 0$）。

(c) $\norm{Tx}^2 = \sum_ns_n^2\abs{\langle e_n, x\rangle}^2
\leq (\max s_n)^2\norm x^2$，等号在使其最大的
$e_n$ 处取得：$\vertiii T = \max s_n$。将奇异值分解在秩
$N$ 处截断，得到范数为 $\sup_{n>N}s_n \to 0$ 的算子：
有限秩逼近。沃尔泰拉算子没有
特征值（见\cref{exo:b3:spectral:7}），但 $V^*V$ 有：
$V^*Vf(x) = \int_x^1\int_0^tf(s)\,\dd s\,\dd t$ 是
对称正核算子（核为 $1 - \max(x,y)$，
弦型格林核），其特征对 --- 借助
边值问题 $-u'' = \lambda^{-1}u$，$u'(0)
= u(1) = 0$（即 \cref{exo:b3:spectral:9} 的族）算出 ---
给出奇异值 $s_n = \bigl((n + \frac12)\pi\bigr)^{-1}$。
奇异值分解之所以依赖于 \emph{两组} 标准正交族，
因为 $V$ 将其特征几何“旋走”：没有特征向量，
却在两组不同的基之间具有完美的对角结构。
\end{solution}

\begin{solution}{pb:b3:spectral:1}
\textbf{1.} 连续性：$\min$ 与 $\max$ 连续；
对称性：交换 $x, y$ 既不改变 $\min(x,y)$，也不改变 $1 -
\max(x,y)$。界：$0 \leq g$，且 $g(x,y) \leq
\max\cdot(1-\max)$ 型的界给出 $g \leq \frac14$（对 $u
= \max$：$\min \leq u$，故 $g \leq u(1-u) \leq \frac14$）。$g
\in L^2(\square)$：是希尔伯特–施密特算子，故 $G$ 紧
（见\cref{exo:b3:spectral:4}）；核是实对称的：$G$
自伴。

\textbf{2.} 在 $y = x$ 处分段：
\[
u(x) = (1 - x)\int_0^x y\,f(y)\,\dd y +
x\int_x^1(1 - y)\,f(y)\,\dd y .
\]
设 $f$ 连续，对上式求导（用乘积法则与
微积分基本定理）：
\[
u'(x) = -\int_0^xyf + \int_x^1(1-y)f
\qquad\text{（边界项相消）},
\]
且 $u''(x) = -xf(x) - (1 - x)f(x) = -f(x)$；显然 $u(0) =
u(1) = 0$。反之，若 $u \in \mathcal C^2$ 在两端取零，
则 $w = u - G(-u'')$ 满足 $w'' = 0$，$w(0) = w(1) =
0$：$w$ 为仿射函数且有两个零点，故 $w = 0$。

\textbf{3.} 设 $Gf = 0$，$f \in L^2$。对 $\psi \in \mathcal
C_c^\infty(\intoo01)$：由第 2 问 $\psi = G(-\psi'')$，于是
\[
\langle f, \psi\rangle = \langle f, G(-\psi'')\rangle
= \langle Gf, -\psi''\rangle = 0
\]
（$G$ 自伴）：由基本引理
（见\cref{cor:b3:lp:fundlemma}），$f = 0$ 几乎处处成立。正性：对
连续 $f$，记 $u = Gf$，
\[
\langle f, Gf\rangle = \int_0^1 fu = \int_0^1(-u'')u
= \bigl[-u'u\bigr]_0^1 + \int_0^1(u')^2 = \int_0^1(u')^2
\geq 0 ;
\]
对 $f \in L^2$，在 $L^2$ 中用连续函数 $f_n$ 逼近：
两端取极限（$G$ 有界）即可。

\textbf{4.} 若 $Gu = \lambda u$，$\lambda \neq 0$：$Gu$
连续（$\abs{Gu(x) - Gu(x')} \leq \norm{g(x,\cdot) -
g(x',\cdot)}_2\norm u_2$，核一致
连续），故 $u$ 有连续代表元；
再由第 2 问的公式知 $Gu \in \mathcal C^2$，于是 $u =
\frac1\lambda Gu \in \mathcal C^2$，且 $-\lambda u'' =
-(Gu)'' = u$，$u(0) = u(1) = 0$。

\textbf{5.} 由 $-\lambda u'' = u$，$u(0) = 0$，得 $u = A\sin(x/\sqrt
\lambda)$（$\lambda$ 为正：由第 3 问，$\lambda =
\langle u, Gu\rangle/\norm u^2 > 0$，对特征向量成立）。$u(1) =
0$ 迫使 $\frac1{\sqrt\lambda} = n\pi$，于是 $\lambda_n =
\frac1{n^2\pi^2}$，特征函数为 $\sin(n\pi x)$，归一化后
$e_n = \sqrt2\sin(n\pi x)$（$\int_0^12\sin^2(n\pi x)\dd x =
1$）。正交性验证：$2\sin(m\pi x)\sin(n\pi x) =
\cos((m-n)\pi x) - \cos((m+n)\pi x)$ 的积分为 $0$，只要 $m
\neq n$。

\textbf{6.} $\ker G = \{0\}$（第 3 问），故
由\cref{thm:b3:spectral:spectral}(1) 得 $H =
\overline{\operatorname{Vect}}(e_n)$：正弦函数构成
$L^2(\intcc01)$ 的希尔伯特基。对 $f(x) = x(1 - x)$：
\[
c_n = \sqrt2\int_0^1x(1-x)\sin(n\pi x)\,\dd x
= \sqrt2\;\frac{2\bigl(1 - (-1)^n\bigr)}{n^3\pi^3}
= \begin{cases}\dfrac{4\sqrt2}{n^3\pi^3} & n \text{ 奇数},\\
0 & n \text{ 偶数},\end{cases}
\]
（两次分部积分）。由帕塞瓦尔恒等式：$\int_0^1x^2(1-x)^2\dd
x = \frac1{30} = \sum_{n \text{ 奇数}}\frac{32}{n^6\pi^6}$，
即 $\sum_{n\text{ 奇数}}n^{-6} = \frac{\pi^6}{960}$。

\textbf{7.} $\langle e_n, Ge_n\rangle = \lambda_n\norm{e_n}^2
= \lambda_n$。固定 $x$，$g(x,
\cdot)$：$\langle e_n, g(x,\cdot)\rangle = (Ge_n)(x) =
\lambda_ne_n(x)$，故 $g(x,\cdot) =
\sum_n\lambda_ne_n(x)\,e_n$ 在 $L^2$ 中成立。显式级数
$\sum_n\lambda_ne_n(x)e_n(y) = \sum_n\frac{2\sin(n\pi
x)\sin(n\pi y)}{n^2\pi^2}$ 在正方形上一致绝对收敛
（$\abs{\text{项}} \leq \frac2{n^2\pi^2}$）：其和
连续，且对每个 $x$，它在
$L^2(\dd y)$ 意义下与 $g(x, \cdot)$ 有相同的系数：两个连续
函数对所有 $(x, y)$ 都相等。令 $y = x$，
并积分（一致绝对收敛允许
逐项积分）：
\[
\int_0^1g(x,x)\,\dd x =
\sum_n\lambda_n\int_0^1e_n(x)^2\dd x = \sum_n\lambda_n .
\]

\textbf{8.} $\int_0^1g(x,x)\dd x = \int_0^1x(1 - x)\dd x =
\frac16$，于是 $\sum_{n\geq1}\frac1{n^2\pi^2} = \frac16$：
\[
\zeta(2) = \sum_{n\geq1}\frac1{n^2} = \frac{\pi^2}6 .
\]

\textbf{9.} 当 $f = \mathbf 1$ 时，$c_n =
\sqrt2\int_0^1\sin(n\pi x)\dd x = \sqrt2\,\frac{1 -
(-1)^n}{n\pi}$：$c_n = \frac{2\sqrt2}{n\pi}$（$n$ 为奇数），偶数时为
$0$。帕塞瓦尔恒等式：$1 = \sum_{n\text{
奇数}}\frac{8}{n^2\pi^2}$，于是 $\sum_{n\text{ 奇数}}n^{-2} =
\frac{\pi^2}8$，且 $\zeta(2) = \frac{\pi^2}8\cdot\frac{1}{1
- \frac14} = \frac{\pi^2}6$（偶数项之和为 $\frac14\zeta(2)$）。
比较：对单个 $f$，帕塞瓦尔恒等式求的是
$\abs{\langle e_n, f\rangle}^2$ 的和；迹公式则
对\emph{核}的对角线取积分，这
相当于对整个标准正交族一次性求帕塞瓦尔和 ---
$\sum_n\langle e_n, Ge_n\rangle$ --- 因而与
任何特定检验函数的选取无关。

\textbf{10.} $\abs{c_n\cos(n\pi t)} \leq \abs{c_n}$，且
$\sum\abs{c_n}^2 < \infty$：对每个 $t$，该级数
在 $L^2$ 中收敛（标准正交展开，
\cref{thm:b3:hilbert:parseval}(3)）；尾项估计
$\norm{u(t) - u_N(t)}_2^2 \leq \sum_{n>N}\abs{c_n}^2$ 对
$t$ 一致成立，且每个部分和关于 $t$ 连续
（只含有限多个余弦）：$t \mapsto u(t,\cdot)$ 是
到 $L^2$ 的连续映射。对于 $f = \sum_{n\leq N}c_ne_n$：每个模式
$\cos(n\pi t)\sin(n\pi x)$ 满足 $\partial_t^2 =
-n^2\pi^2 = \partial_x^2$ 作用于其上，在 $x = 0,
1$ 处为零，取值 $\sin(n\pi x)$，时间导数为 $0$，二者皆在 $t =
0$ 时成立：有限和即为全部解。从音乐角度看，弦的
运动是\emph{驻波} $e_n$ 的叠加，
其频率 $n\pi$ 为基频及其
泛音；谱定理断言，每个初始形状
唯一分解为这些纯音，系数
$c_n$ 即音色。听弦声，就是在计算一个
标准正交展开。

\textbf{11.} 记 $a_n = \abs{\langle u_n, x\rangle}^2$：
$m_{p+1} = \sum_n\mu_n^{p+1}a_n = \sum_n\bigl(\mu_n^{p/2}
\sqrt{a_n}\bigr)\bigl(\mu_n^{p/2+1}\sqrt{a_n}\bigr) \leq
\sqrt{m_p\,m_{p+2}}$（$\ell^2$ 中的柯西–施瓦茨不等式）。由此
比值 $m_{p+1}/m_p$ 关于 $p$ 单调不减；由于
$R(x) = \frac{m_1}{m_0}$，$\frac{\langle Ax, Ax\rangle}
{\langle x, Ax\rangle} = \frac{m_2}{m_1}$ 及 $R(Ax) =
\frac{m_3}{m_2}$，链式不等式由此成立，每一项 $\leq \mu_1$
因为 $m_{p+1} \leq \mu_1m_p$ 逐项成立。收敛性：若
$a_1 > 0$（将最大特征值的总权重记为
$a_1$），则
\[
\mu_1 \geq R(A^kx) = \frac{m_{2k+1}}{m_{2k}} =
\mu_1\,\frac{a_1 + \sum_{\mu_n<\mu_1}(\mu_n/\mu_1)^{2k+1}
a_n}{a_1 + \sum_{\mu_n<\mu_1}(\mu_n/\mu_1)^{2k}a_n}
\longrightarrow \mu_1,
\]
由和式的控制收敛（比值 $< 1$）：
幂法对每个不与最大特征子空间
正交的初始向量都收敛。

\textbf{12.} 逐点地，$\abs{\langle x, (A - B)x\rangle}
\leq \vertiii{A - B}\,\norm x^2$，故 $R_A(x) \leq R_B(x) +
\vertiii{A - B}$ 对每个 $x$ 成立。将此代入
\cref{exo:b3:spectral:8}的极大–极小公式：$\mu_n(A)
\leq \mu_n(B) + \vertiii{A - B}$，并关于 $A,
B$ 对称地得到 $\abs{\mu_n(A) - \mu_n(B)} \leq \vertiii{A - B}$ 对
所有 $n$ 同时成立。

\textbf{13.} $-w'' = x - x^2$ 积分得 $w = -\frac{x^3}6
+ \frac{x^4}{12} + cx$（$w(0) = 0$），由 $w(1) = 0$
得 $c = \frac1{12}$：
\[
w = \frac{x^4 - 2x^3 + x}{12} = \frac{x(1-x)(1 + x -
x^2)}{12} = Gu .
\]
于是 $\norm u_2^2 = \int_0^1x^2(1-x)^2 = \frac1{30}$，
并利用 $\int_0^1x^3(1-x)^3 = B(4,4) = \frac1{140}$：
\[
\langle u, Gu\rangle = \frac1{12}\Bigl(\frac1{30} +
\frac1{140}\Bigr) = \frac{17}{5040},
\qquad
R(u) = \frac{17/5040}{1/30} = \frac{17}{168} .
\]
于是 $\frac1{\pi^2} = \lambda_1 \geq \frac{17}{168}$，即
$\pi^2 \leq \frac{168}{17} = 9.8824$：$\pi \leq 3.14364 <
3.1437$（真值为 $\pi^2 = 9.8696$）。一个多项式、一个
积分、一个数字。

\textbf{14.} 记 $u = x - x^2$，故 $Gu =
\frac{u(1 + u)}{12}$，并利用 $\int u^2 = \frac1{30}$，
$\int u^3 = \frac1{140}$，$\int u^4 = B(5,5) =
\frac{4!\,4!}{9!} = \frac1{630}$：
\[
\norm{Gu}_2^2 = \frac1{144}\int u^2(1+u)^2 =
\frac1{144}\Bigl(\frac1{30} + \frac2{140} +
\frac1{630}\Bigr) = \frac1{144}\cdot\frac{62}{1260} =
\frac{31}{90720} .
\]
于是 $\frac{\langle Gu, Gu\rangle}{\langle u, Gu\rangle} =
\frac{31/90720}{17/5040} = \frac{31}{306}$，并由第
11 题，此值仍 $\leq \mu_1 = \frac1{\pi^2}$：$\pi^2 \leq
\frac{306}{31} = 9.87097$，即 $\pi \leq 3.14181 <
3.1419$ --- 四位数字（而下一次迭代给出
约八位，误差每步以
$(\lambda_2/\lambda_1)^2 = \frac1{16}$ 为因子收缩）。

\textbf{15.} 族 $(e_m \otimes e_n)(x,y) =
e_m(x)e_n(y)$ 是 $L^2(\intcc01^2)$ 的希尔伯特基
（标准正交性由托内利定理得出；完备性与
\cref{exo:b3:spectral:5}中的相同）。由第 7 题，对固定的 $x$：
$g(x, \cdot) = \sum_n\lambda_ne_n(x)e_n$，故
$g$ 在 $e_m\otimes e_n$ 上的系数为
\[
\langle e_m\otimes e_n,\ g\rangle
= \int_0^1 e_m(x)\,\lambda_n e_n(x)\,\dd x
= \lambda_n\,\delta_{mn} .
\]
正方形上的帕塞瓦尔恒等式：
\[
\iint g^2 = \sum_{m,n}\abs{\langle e_m\otimes e_n,
g\rangle}^2 = \sum_n\lambda_n^2 .
\]

\textbf{16.} 由 $g$ 的对称性，
\[
\iint g^2 = 2\iint_{y<x}y^2(1-x)^2
= 2\int_0^1(1-x)^2\,\frac{x^3}3\,\dd x
= \frac23\,B(4, 3) = \frac23\cdot\frac{3!\,2!}{6!}
= \frac23\cdot\frac1{60} = \frac1{90} .
\]

\textbf{17.} 合并可得 $\sum_n\frac1{n^4\pi^4} =
\frac1{90}$，即 $\zeta(4) = \frac{\pi^4}{90}$。
一般地，$\operatorname{tr}(G^k) = \sum\lambda_n^k =
\frac{\zeta(2k)}{\pi^{2k}}$ 等于由有理多项式核 $g$ 的若干乘积
所作的累次积分，积分区域为
$k$ 维立方体，其值为有理数。于是 $\zeta(2k) \in
\pi^{2k}\Q$，对一切 $k$ 成立。该方法只能给出偶数
自变量的结果，原因在于特征值是通过其
\emph{幂} --- $\sum\lambda_n^k$ --- 出现的，而 $\lambda_n =
\frac1{n^2\pi^2}$：任何迹的组合都无法得出 $\sum
n^{-3}$；$\zeta(3)$ 的算术性质（由阿佩里证明为无理数，
超越性仍未解决），已超出谱的
计算范围。

\textbf{18.} 正项之和中的最大项至多
不超过整个和：$\lambda_1^2 \leq \sum\lambda_n^2 =
\frac1{90}$，于是 $\frac1{\pi^2} \leq \frac1{\sqrt{90}}$，从而
$\pi \geq 90^{1/4} = 3.0801\ldots$。结合第 14 题，
$3.080 < \pi < 3.1419$，仅凭弦的运算即可。更高的
迹以几何速率收紧下界：
$\lambda_1 \leq (\operatorname{tr}G^{2k})^{1/2k} =
\lambda_1\bigl(1 + \sum_{n\geq2}(\lambda_n/\lambda_1)^{2k}
\bigr)^{1/2k}$，而寄生因子以
$\bigl(\tfrac14\bigr)^{2k}\cdot\frac1{2k}$ 的速度趋于零 --- 与
幂法收敛所依赖的谱间隙机制相同，
（第 11 题），只是此处从迹的一侧来看。

\textbf{19.} $\abs{n^2\pi^2 - \nu} \geq \delta > 0$ 对所有
$n$ 成立（数列 $n^2\pi^2 \to \infty$ 以一定的
余量避开 $\nu$），且 $\abs{n^2\pi^2 - \nu} \geq \frac{n^2\pi^2}2$
对充分大的 $n$ 时成立。$L^2$ 收敛：系数
$\frac{c_n}{n^2\pi^2 - \nu}$ 平方可和（
被 $\frac{\abs{c_n}}\delta$ 控制）。一致收敛：尾部的
上确界范数以下式为界：$\sqrt2\sum_{n>N}
\frac{\abs{c_n}}{\abs{n^2\pi^2 - \nu}} \leq
\frac{2\sqrt2}{\pi^2}\bigl(\sum\abs{c_n}^2\bigr)^{1/2}
\bigl(\sum_{n>N}n^{-4}\bigr)^{1/2} \to 0$（柯西–施瓦茨）。
验证：$Gf + \nu Gu$ 的 $e_n$ 系数为
\[
\lambda_nc_n + \frac{\nu\lambda_nc_n}{n^2\pi^2 - \nu}
= \frac{c_n}{n^2\pi^2}\Bigl(1 + \frac{\nu}{n^2\pi^2 -
\nu}\Bigr) = \frac{c_n}{n^2\pi^2 - \nu} :
\]
恰为
$u$ 的系数，于是 $u = Gf + \nu Gu$；唯一性
可由两解之差 $v$ 满足 $v = \nu
Gv$，即 $\langle e_n, v\rangle(n^2\pi^2 - \nu) = 0$，对
所有 $n$ 成立，故 $v = 0$。

\textbf{20.} 与第 19 题相同，方程 $u = Gf + \nu
Gu$ 等价于下列系数方程族
$(n^2\pi^2 - \nu)\,\langle e_n, u\rangle = c_n$，$n \geq
1$。对 $n = m$，左边为
$0$：可解性要求 $c_m = 0$，于是 $\langle e_m,
u\rangle$ 可任意取值，其余系数均
被确定：解集构成直线 $u_0 + \R e_m$。
共振：若外力在某特征模上有分量，
则会不断向其输入能量，使之无界增长 --- 恰如在秋千的
固有频率上推动秋千。

\textbf{21.} 由第 19 题的公式，$\norm{R_\nu f}_2^2
= \sum\frac{\abs{c_n}^2}{(n^2\pi^2 - \nu)^2} \leq
\frac{\norm f^2}{(\pi^2 - \nu)^2}$（当 $\nu < \pi^2$ 时，
最近的特征值为 $\pi^2$），等号可在
$f = e_1$ 处逼近，故算子范数为 $\frac1{\pi^2 - \nu}$。紧性：
$R_\nu$ 是其有限秩截断的范数极限
（尾部系数 $\frac1{n^2\pi^2 - \nu} \to 0$）；
自伴性与正性可由对角
形式读出（所有系数 $\frac1{n^2\pi^2 - \nu} > 0$）。
$R_\nu$ 的特征值为 $\frac1{n^2\pi^2 - \nu}$，
I--VI 各部分的分析可原样照搬。

\textbf{22.} 对照如下：\emph{特征值} $\lambda_n =
\frac1{n^2\pi^2}$ $\leftrightarrow$ 平方频率
$n^2\pi^2$（第 $n$ 个谐波）；\emph{迹} $\sum
\lambda_n = \frac16$ $\leftrightarrow$ $\zeta(2) =
\frac{\pi^2}6$；\emph{希尔伯特–施密特范数} $\iint g^2 =
\frac1{90}$ $\leftrightarrow$ $\zeta(4) = \frac{\pi^4}{90}$；
\emph{弗雷德霍姆择一性} $\leftrightarrow$ 受迫弦的
共振；\emph{极小极大} $\leftrightarrow$
变分估计，由一个多项式即可得到 $\pi < 3.1437$ 的
界。每一组对应的背后，都是同一个对象：一个
紧自伴算子，对角化一次，即可从
五个方面加以利用。

\textbf{23.} 按奇偶性拆分，并在偶数部分中以 $n = 2m$
代入，得
\[
\zeta(6) = \sum_{n\text{ 奇数}}\frac1{n^6} +
\sum_{m\geq1}\frac1{(2m)^6}
= \frac{\pi^6}{960} + \frac{\zeta(6)}{64},
\]
故 $\frac{63}{64}\zeta(6) = \frac{\pi^6}{960}$，且 $\zeta(6)
= \frac{64\,\pi^6}{63\cdot960} = \frac{\pi^6}{945}$。若改用迹，
则需计算 $\operatorname{tr}G^3 =
\sum\lambda_n^3 = \zeta(6)/\pi^6$ 写成 $\iint g\,g_2$，其中迭代核
为 $g_2(x,y) = \int_0^1g(x,z)g(z,y)\dd z$ ---
需对分段多项式作三次积分；帕塞瓦尔
恒等式用于 $x(1-x)$ 时仅需一次。

\textbf{24.} (a) 对角化：$v = \sum_na_nu_n$（可附加一个
核分量，对其 $\langle v, Av\rangle$
无增益，而 $\norm v^2$ 增大，故极大元
不含该分量）。于是 $\langle v, Av\rangle = \sum\mu_na_n^2 \leq
\mu_1\sum a_n^2$，等号成立当且仅当 $a_n = 0$
当 $\mu_n < \mu_1$ 时，极大元位于
属于 $\mu_1$ 的特征子空间。(b) 对任意 $u$，
\[
\langle\abs u, A\abs u\rangle - \langle u, Au\rangle
= \iint k(x,y)\,\bigl(\abs{u(x)}\abs{u(y)} -
u(x)u(y)\bigr)\dd x\,\dd y \;\geq\; 0,
\]
被积函数逐点非负。若 $P = \{u >
0\}$ 与 $N = \{u < 0\}$ 的测度均为正，则在
$P\times N$ 上的被积函数等于 $2k\abs{u(x)}\abs{u(y)} >
0$ 于一正测度集上，严格不等式成立。
属于 $\mu_1$ 的特征函数 $u$ 使瑞利商取到最大值，
而 $\abs u$ 的范数相同，故严格性将给出
$R(\abs u) > \mu_1$ --- 不可能；因此 $u$ 几乎处处
具有定号，不妨设 $u \geq 0$。于是 $u(x) = \mu_1^{-1}(Au)(x) =
\mu_1^{-1}\int k(x,y)u(y)\dd y > 0$ 对每个 $x \in
\intoo01$ 成立（$k(x,\cdot) > 0$ 且 $u \neq 0$）。(c) 若
特征子空间的维数 $\geq 2$，则其中含两个
\emph{正交}特征函数 $u, v$，每个都具有定号，
且由 (b) 知其在内部处处不为零；然而此时
$\abs{\langle u, v\rangle} = \int\abs u\,\abs v > 0$ ---
矛盾。就弦而言：$k = g > 0$ 在开
正方形上，$\mu_1 = \lambda_1 = \frac1{\pi^2}$ 确为
单特征值，$e_1 = \sqrt2\sin(\pi x)$ 在
$\intoo01$ 上为正；而每个 $e_n = \sqrt2\sin(n\pi x)$，$n \geq 2$，
与正的 $e_1$ 正交，故与之乘积的积分必为零，
因而必变号 --- 正如其 $n - 1$ 个内部
零点 $\frac kn$ 所印证。

\textbf{25.} 由于 $n^2\pi^2 \to \infty$，最小值 $d =
\min_n\abs{n^2\pi^2 - \nu}$ 可取到，设其在某模态 $m$ 处取得，
且 $d > 0$，因为 $\nu$ 避开了谱。第 19 题的
对角公式给出
\[
\norm{R_\nu f}_2^2 =
\sum_n\frac{\abs{c_n}^2}{(n^2\pi^2 - \nu)^2}
\leq \frac1{d^2}\,\norm f_2^2,
\]
等号在 $f = e_m$ 时取得：$\vertiii{R_\nu} = \frac1d$，
即 $\nu$ 到谱的距离之倒数，
--- 这是一般的预解式原理，此处以显式
坐标给出。自伴性可由实的对角
系数直接读出；紧性与第 21 题相同（
各系数趋于 $0$，故有限秩截断
在范数意义下收敛）。共振代价：对 $f = e_1$，取 $\nu =
(1-\varepsilon)\pi^2$，公式给出 $u =
\frac{c_1}{\pi^2 - \nu}e_1 = \frac1{\varepsilon\pi^2}e_1$,
与静态响应 $Ge_1 = \frac1{\pi^2}e_1$ 相比：
放大倍数为 $\frac1\varepsilon$。取 $\varepsilon =
10^{-2}$ 时，响应为静态响应的 $100$ 倍 --- 且
它随 $\varepsilon \to 0$ 而发散，这正是第 20 题的择一性
从有界一侧看到的体现：强迫频率越
接近某个固有频率，逆
算子便越不有界。
\end{solution}
